% Generated draft: requires mathematical and source review.
\documentclass[11pt]{article}
\usepackage[margin=25mm]{geometry}
\usepackage{fontspec}
\setmainfont{Latin Modern Roman}
\usepackage{amsmath,amssymb,mathtools,unicode-math}
\setmathfont{Latin Modern Math}
\usepackage{enumitem,needspace,xurl,hyperref}
\setlength{\parindent}{0pt}
\setlength{\parskip}{5pt}
\setlength{\emergencystretch}{3em}
\newcommand{\N}{\mathbb{N}}
\newcommand{\Z}{\mathbb{Z}}
\newcommand{\Q}{\mathbb{Q}}
\newcommand{\R}{\mathbb{R}}
\newcommand{\C}{\mathbb{C}}
\newcommand{\F}{\mathbb{F}}
\newcommand{\A}{\mathbb{A}}
\newcommand{\PP}{\mathbb{P}}
\newcommand{\E}{\mathbb{E}}
\newcommand{\eps}{\varepsilon}
\newcommand{\dd}{\,\mathrm{d}}
\begin{document}
\section*{UnsolvedMath ID: 30003708}
\textbf{Source problem code:} catalogue entry
\subsection{Definitions and mathematical statement}
The supplied question concerns a field \(k\) described as a ``maximal ordered field.'' The record does not define maximality or specify the extensions relative to which it is taken; this hypothesis is therefore retained as an unresolved source qualification.

Under the conventional signature interpretation of the orthogonal-group notation, define, for nonnegative integers \(n,i\),
\[
q_{n,i}(x,y)=\sum_{a=1}^{n}x_a^2-\sum_{b=1}^{i}y_b^2,
\qquad (x,y)\in k^n\oplus k^i,
\]
\[
O_{n,i}(k)=\{g\in\mathrm{GL}_{n+i}(k):q_{n,i}(gv)=q_{n,i}(v)\text{ for every }v\in k^{n+i}\},
\qquad O_m(k)=O_{m,0}(k).
\]
The negative definite form \(-\sum y_b^2\) has the same isometry group as \(\sum y_b^2\), so block diagonal matrices give an inclusion
\[
\iota_{n,i}:O_n(k)\times O_i(k)\hookrightarrow O_{n,i}(k),
\qquad (A,B)\longmapsto\operatorname{diag}(A,B).
\]
Here \(H_d(G;\Z)\) denotes homology of the abstract group \(G\), regarded as discrete, with trivial integral coefficients. If zero-dimensional forms are included, their orthogonal groups are trivial.

\textbf{Question, conditional on these conventions.} For every field \(k\) satisfying the intended maximal-ordered-field hypothesis and every fixed pair of nonnegative integers \(d,i\), does there exist a nonnegative integer \(N=N(k,d,i)\) such that, for every integer \(n\geq N\), the map
\[
(\iota_{n,i})_*:
H_d\bigl(O_n(k)\times O_i(k);\Z\bigr)
\longrightarrow H_d\bigl(O_{n,i}(k);\Z\bigr)
\]
is an isomorphism?

The form conventions and nonnegative indexing above make explicit a conventional reading, but are not defined in the supplied source record. The wording ``sufficiently large relative to \(d\) and \(i\)'' may additionally require a bound \(N(d,i)\) uniform over the permitted fields \(k\); the record does not resolve this quantifier distinction. No restriction on \(i\) is supplied. These qualifications require source clarification before the displayed formulation can be treated as definitive.
\subsection{Short English statement}
For a field satisfying the source's undefined ``maximal ordered'' hypothesis, ask whether the block diagonal inclusion of two definite orthogonal groups into the corresponding indefinite orthogonal group induces an isomorphism on integral discrete group homology once \(n\) is sufficiently large for fixed \(d\) and \(i\). The meaning of maximality, the precise form and index conventions, and uniformity in the field remain unresolved.
\subsection{Sources}
\begin{itemize}
\item Local registry: \texttt{data/unsolvedmath/problems.json}, with \texttt{catalog/problems.json} as the additional-record registry; supplied permanent UnsolvedMath ID \(30003708\), source problem code \texttt{catalogue entry}, title \emph{Homological Stability for Indefinite Orthogonal Groups}. No individual proposer is supplied. The recorded proposed year is 2018.
\item Supplied report citation: \emph{Topology of Arrangements and Representation Stability (2019)}, Oberwolfach Reports, DOI \texttt{10.4171/owr/2018/2}, pp.\ 43--123. \url{https://doi.org/10.4171/owr/2018/2}. The background attributes the problem collection to workshop open-problem discussions documented in Oberwolfach Reports. The supplied title year 2019 differs from the report designation 2/2018 and the proposed year; this discrepancy is unresolved.
\item Supplied literature reference: Sprehn--Wahl, ``classical-group stability.'' \url{https://arxiv.org/abs/1812.08742}.
\item Supplied literature reference: Schlichting--Sood, ``orthogonal stability.'' \url{https://arxiv.org/abs/2206.00981}.
\item Supplied literature reference: Oberwolfach Report 2/2018. \url{https://doi.org/10.4171/owr/2018/2}.
\item Literature assessment dated \texttt{2026-08-21}: the supplied curation labels the problem partially solved and states that homological stability is proved for broad families of orthogonal groups, while those results do not completely settle this exact mixed-form inclusion over a maximal ordered field for the full stated range of \(i\). This is supplied dated curation derived from the OWR material, not an independently verified assessment. The record also labels the statement \texttt{corrected\_verified}, explaining that its mathematical content was preserved while notation, typography, or omitted definitions were repaired; this supplied label does not establish mathematical correctness of this draft.
\end{itemize}
\end{document}
